%titolo{C'è qualcuno che mi può aiutare}
%autore{}
%album{}
%tonalita{DO}
%famiglia{Scout}
%gruppo{}
%momenti{}
%identificatore{ce_qualcuno_che_mi_puo_aiutare}
%data_revisione{2026_07_1}
%trascrittore{Lunardelli Luca}

\beginsong{C'è qualcuno che mi può aiutare}

\beginverse
\[C]Bosco silen\[F]zioso \[G]Cocci cammina qua e \[C]là 
\[C7]una strana foglia a for\[F]ma di esse 
\[A-]ripete sibilando questa fra\[C]se \[G]qua
\endverse

\beginchorus
\[C]C'è qualcuno che mi \[F]può aiutare? 
C'è qual\[G]cuno! che! mi! può! \[F]aiutare? 
\[C]C'È QUALCUNO CHE MI \[F]PUÒ AIUTARE!
\[G]SE MI SEN\[C]TIRÀ 
\[C]dimmi cosa posso fa\[F]re per te 
\[G]e lo farò volent\[C]ie\[B]ri \[G]a te 
\[C]dimmi cosa posso fa\[F]re per te 
\[G]e lo farò volen\[C]tier \rep{2}
\endchorus
\chordsoff
\beginverse
Vento burrascoso ha portato una scheggia qua si è infilata
nel mio occhio e vorrei che uscisse di là
\endverse

\beginchorus
C'è qualcuno che mi può aiutare? 
C'è qualcuno! che! mi! può! aiutare?
C'È QUALCUNO CHE MI PUÒ AIUTARE!
SE MI SENTIRÀ 
Ecco cosa posso fare per te 
e lo farò volentieri a te 
Ecco cosa posso fare per te 
e lo farò volentier \rep{2}
\endchorus

\beginverse
Cocci amica coccinella mi hai guarito ti ringrazio già
sono a tua disposizione in groppa ti porta Scibà
\endverse

\beginchorus
C'è qualcuno che mi può aiutare? 
C'è qualcuno! che! mi! può! aiutare? 
C'È QUALCUNO CHE MI PUÒ AIUTARE!
SE MI SENTIRÀ
Ecco cosa posso fare per te 
e lo farò volentieri a te
Ecco cosa posso fare per te 
e lo farò volentier \rep{2}
\endchorus
\endsong
